\chapter{Vector bundles}
\label{chap:bundles}

Vector bundles with fibres modeled on locally convex spaces, defined by trivializing covers,
pullbacks, Hom and dual bundles and the cotangent bundle, sections, connections on vector bundles
with their curvature, induced connections and bundle metrics. As for connections on the tangent
bundle (Chapter~\ref{chap:connections}), bilinear families are encoded as maps of several
variables, with joint smoothness.

\section{Vector bundles and their connections}
\label{sec:vector-bundles}

\begin{definition}[Trivializing covers]
  \label{def:trivialization-cover}
  \lean{LocallyConvex.Manifolds.TrivializationCover,
    LocallyConvex.Manifolds.TrivializationCover.domain,
    LocallyConvex.Manifolds.TrivializationCover.transDomain,
    LocallyConvex.Manifolds.TrivializationCover.trans,
    LocallyConvex.Manifolds.TrivializationCover.smoothMBOn_trans,
    LocallyConvex.Manifolds.atlasCover}
  \leanok
  \uses{def:mb-manifold}
  A \emph{trivializing cover} of $M$, indexed by a set $\iota$, is a family of open sets
  $U_i \subseteq M$, each contained in the domain of a chart $\varphi_i$ of the atlas, together
  with a choice, for every $x \in M$, of an index $i(x)$ with $x \in U_{i(x)}$. The domains of
  the charts of the atlas form one.
\end{definition}

\begin{definition}[Vector bundles]
  \label{def:vector-bundle}
  \lean{LocallyConvex.Manifolds.MBVectorBundleCore,
    LocallyConvex.Manifolds.MBVectorBundleCore.coordChangeLinear,
    LocallyConvex.Manifolds.MBVectorBundleCore.continuousOn_coordChange,
    LocallyConvex.Manifolds.MBVectorBundleCore.toFiberBundleCore,
    LocallyConvex.Manifolds.MBVectorBundleCore.TotalSpace,
    LocallyConvex.Manifolds.MBVectorBundleCore.bundleChart}
  \leanok
  \uses{def:trivialization-cover, def:tangent-bundle}
  A \emph{vector bundle} over $M$ with fiber $F$, trivialized over a trivializing cover
  $(U_i, \varphi_i)$, is given by coordinate changes $g_{ij}(x) \colon F \to F$ between the
  trivializations over $U_i$ and $U_j$, linear for $x \in U_i \cap U_j$, satisfying the cocycle
  conditions, and such that $(y, \xi) \mapsto g_{ij}(\varphi_i^{-1} y)\xi$ is smooth on
  $\varphi_i(U_i \cap U_j) \times F$.
\end{definition}

If $F$ is not normable, $L(F)$ carries no vector topology for which the evaluation is continuous
\cite[Satz 1]{Mai63},
so one cannot ask $y \mapsto g_{ij}(y)$ to be a smooth map into $GL(F)$; joint smoothness of
$(y, \xi) \mapsto g_{ij}(y)\xi$ is the substitute. For Banach spaces it is the classical
condition only up to this difference of topologies on $L(F)$.

\begin{theorem}[The total space is a manifold]
  \label{thm:vector-bundle-manifold}
  \lean{LocallyConvex.Manifolds.MBVectorBundleCore.smoothMBOn_bundleChart_transition,
    LocallyConvex.Manifolds.MBVectorBundleCore.bundleChart_transition_mem,
    LocallyConvex.Manifolds.MBVectorBundleCore.mbManifold,
    LocallyConvex.Manifolds.tangentMBCore}
  \leanok
  \uses{def:vector-bundle}
  The total space of a vector bundle with fiber $F$ is a manifold modeled on $E \times F$, with
  transition maps $(y, \xi) \mapsto (\tau y, g(y)\xi)$. The tangent bundle is a vector bundle,
  with coordinate changes $(y, w) \mapsto \dd\tau(y) w$.
\end{theorem}

\begin{proof}
  \leanok
  As for the tangent bundle: the transition maps are compositions of transition maps of $M$ with
  the smooth maps $(y, \xi) \mapsto g(y)\xi$.
\end{proof}

\begin{definition}[Pullback bundles]
  \label{def:pullback-bundle}
  \lean{LocallyConvex.Manifolds.TrivializationCover.pullback,
    LocallyConvex.Manifolds.SmoothMBM.continuous,
    LocallyConvex.Manifolds.MBVectorBundleCore.pullback,
    LocallyConvex.Manifolds.MBVectorBundleCore.pullbackMap,
    LocallyConvex.Manifolds.MBVectorBundleCore.pullbackMap_proj}
  \leanok
  \uses{def:vector-bundle, def:contmdiff}
  Let $f \colon N \to M$ be smooth and $V$ a vector bundle over $M$, trivialized over
  $(U_i, \varphi_i)$. The \emph{pullback} $f^*V$ is the vector bundle over $N$ trivialized over the
  sets $k.\mathrm{source} \cap f^{-1}(U_i)$, for $k$ a chart of $N$, with coordinate changes
  $g_{ij} \circ f$. Its fiber over $x$ is the fiber of $V$ over $f(x)$.
\end{definition}

The bundle $f^*V$ need not be trivial over the domain of a chart of $N$, since $f$ need not map
it into a single $U_i$; this is why vector bundles are trivialized over general covers rather
than over the atlas of the base.

\begin{theorem}[Hom bundles, dual bundles and the cotangent bundle]
  \label{thm:hom-bundle}
  \lean{LocallyConvex.Manifolds.MBVectorBundleCore.continuous_coordChange,
    LocallyConvex.Manifolds.MBVectorBundleCore.coordCLM,
    LocallyConvex.Manifolds.MBVectorBundleCore.coordCLM_apply,
    LocallyConvex.Manifolds.MBVectorBundleCore.smoothMBOn_coordCLM,
    LocallyConvex.Manifolds.MBVectorBundleCore.smoothMBOn_coordCLM_symm,
    LocallyConvex.Manifolds.homCore, LocallyConvex.Manifolds.homCore_coordChange_apply,
    LocallyConvex.Manifolds.trivialCore, LocallyConvex.Manifolds.dualCore,
    LocallyConvex.Manifolds.cotangentCore}
  \leanok
  \uses{def:vector-bundle, thm:vector-bundle-manifold}
  Let $V_1$, $V_2$ be vector bundles over $M$ with normed fibers $F_1$, $F_2$. There is a vector
  bundle $\mathrm{Hom}(V_1, V_2)$ with fiber $\Lin(F_1, F_2)$ (operator norm) and coordinate
  changes $A \mapsto g_2(x) \circ A \circ g_1(x)^{-1}$. In particular there is a dual bundle
  $V' = \mathrm{Hom}(V, \R)$, and if $E$ is normed, a cotangent bundle $T^*M = (TM)'$, with
  no completeness or reflexivity assumption.
\end{theorem}

\begin{proof}
  \leanok
  The cocycle conditions are those of $V_1$ and $V_2$. For smoothness, the maps
  $y \mapsto g_{ij}(y)$ are smooth for the operator norm, because $(y, \xi) \mapsto g_{ij}(y)\xi$
  is smooth (\cite{LCA}, theorem on operator-valued maps, \texttt{contDiffMBOn\_of\_unc\_clm}),
  and composition of operators is continuous and bilinear.
\end{proof}

If a fiber is not normable, $L(F_1, F_2)$ carries no vector topology for which the evaluation is
continuous (\cite[Satz 1 and its footnote]{Mai63}), and these constructions are not available;
Lempert notes the same restriction
\cite[\S2, p.~488]{Lem98}.

\begin{definition}[Connections on vector bundles]
  \label{def:bundle-connection}
  \lean{LocallyConvex.Manifolds.BundleConnection,
    LocallyConvex.Manifolds.localCovDerivB}
  \leanok
  \uses{def:vector-bundle, def:connection-forms}
  A \emph{connection} on a vector bundle is given by symbols $\Gamma_i(y)(a, \xi) \in F$ in the
  trivialization over each chart, smooth and bilinear, subject to the gauge transformation law
  \[
    \Gamma_j(\tau y)(\dd\tau(y) a, g(y)\xi) = g(y)\Gamma_i(y)(a, \xi) + (\partial_a g)(y)\xi,
  \]
  where $\tau = j \circ i^{-1}$ and $g(y) = g_{ij}(i^{-1} y)$. In a trivialization, the covariant
  derivative of a section is $\nabla_X s = \dd s \cdot X - \Gamma(X, s)$.
\end{definition}

\begin{lemma}[Calculus of the covariant derivative]
  \label{lem:bundle-cov-deriv}
  \lean{LocallyConvex.Manifolds.dMB_fiber_comp,
    LocallyConvex.Manifolds.localCovDerivB_transition,
    LocallyConvex.Manifolds.IsLocalSymbol.dMB_apply,
    LocallyConvex.Manifolds.IsLocalSymbol.localCovDerivB_curv,
    LocallyConvex.Manifolds.IsLocalSymbol.contDiffMBOn_localCovDerivB,
    LocallyConvex.Manifolds.IsLocalSymbol.mono}
  \leanok
  \uses{def:bundle-connection, lem:conn-local-curv}
  Under the gauge transformation law, the covariant derivative transforms covariantly:
  if $X_j(\tau y) = \dd\tau(y) X_i(y)$ and $s_j \circ \tau = g \cdot s_i$, then
  $\nabla_{X_j} s_j(\tau y) = g(y) \nabla_{X_i} s_i(y)$. For $C^2$ fields and sections,
  $\nabla_X\nabla_Y s - \nabla_Y\nabla_X s - \nabla_{[X, Y]} s = \Omega(X, Y) s$ pointwise.
\end{lemma}

\begin{proof}
  \leanok
  Product rules for $y \mapsto g(y)s(y)$ and $y \mapsto \Gamma(y)(A y, B y)$; in the curvature
  formula the terms with derivatives of the fields and the section cancel, using the symmetry of
  second derivatives of $s$, as in Lemma~\ref{lem:conn-local-curv}.
\end{proof}

\begin{theorem}[The curvature is a tensor]
  \label{thm:curvature-gauge}
  \lean{LocallyConvex.Manifolds.curvature_gauge,
    LocallyConvex.Manifolds.BundleConnection.curvLoc_transition,
    LocallyConvex.Manifolds.BundleConnection.dLoc_curvLoc,
    LocallyConvex.Manifolds.Connection.toBundleConnection,
    LocallyConvex.Manifolds.Connection.localCurv_transition}
  \leanok
  \uses{lem:bundle-cov-deriv, thm:second-bianchi, def:connection}
  The curvature forms of a connection on a vector bundle satisfy
  $\Omega_j(\tau y)(\dd\tau\, u, \dd\tau\, v)(g\xi) = g\,\Omega_i(y)(u, v)\xi$, and the second
  Bianchi identity holds in every trivialization. A connection on $M$ is a connection on the
  tangent bundle; in particular its curvature is a tensor.
\end{theorem}

\begin{proof}
  \leanok
  Apply the curvature formula of Lemma~\ref{lem:bundle-cov-deriv} to the constant fields $u$,
  $v$ and the constant section $\xi$ in the chart $i$, and to their images
  $\dd\tau(\sigma z)u$, $\dd\tau(\sigma z)v$, $g(\sigma z)\xi$ in the chart $j$, where
  $\sigma = \tau^{-1}$. Covariant derivatives and brackets transform covariantly, so the
  left-hand sides correspond under $g(y)$. The argument uses neither partitions of unity nor
  bump functions, which need not exist on locally convex manifolds.
\end{proof}

\begin{definition}[Sections and the covariant derivative]
  \label{def:bundle-sections}
  \lean{LocallyConvex.Manifolds.MBVectorBundleCore.IsSection,
    LocallyConvex.Manifolds.MBVectorBundleCore.localRep,
    LocallyConvex.Manifolds.MBVectorBundleCore.bundleChart_mem_atlas,
    LocallyConvex.Manifolds.MBVectorBundleCore.localRep_transition,
    LocallyConvex.Manifolds.MBVectorBundleCore.localRep_indexAt,
    LocallyConvex.Manifolds.MBVectorBundleCore.contDiffMBOn_localRep,
    LocallyConvex.Manifolds.BundleConnection.covDeriv,
    LocallyConvex.Manifolds.BundleConnection.isSection_covDeriv,
    LocallyConvex.Manifolds.BundleConnection.covDeriv_snd,
    LocallyConvex.Manifolds.localCovDerivB_congr_right}
  \leanok
  \uses{def:bundle-connection, def:vector-field}
  A section $s$ of a vector bundle has local representatives $s_i$ with
  $s_j(\tau y) = g(y) s_i(y)$; those of a $C^k$ section are $C^k$. For a connection, the
  covariant derivative $\nabla_X s$ of a section along a vector field is defined in the
  preferred chart at each point by $\dd s_i \cdot X_i - \Gamma_i(X_i, s_i)$.
\end{definition}

\begin{theorem}[The global covariant derivative and curvature]
  \label{thm:bundle-cov-deriv-global}
  \lean{LocallyConvex.Manifolds.BundleConnection.localRep_covDeriv,
    LocallyConvex.Manifolds.BundleConnection.curvature,
    LocallyConvex.Manifolds.BundleConnection.curvature_eq}
  \leanok
  \uses{def:bundle-sections, lem:bundle-cov-deriv}
  For a $C^1$ section, $\nabla_X s$ is given by $\dd s_j \cdot X_j - \Gamma_j(X_j, s_j)$ in
  every trivialization. For $C^2$ fields and sections, the curvature
  $R(X, Y)s = \nabla_X\nabla_Y s - \nabla_Y\nabla_X s - \nabla_{[X, Y]} s$ at $x$ is the
  curvature form of the preferred chart evaluated at the values of $X$, $Y$, $s$ at $x$.
\end{theorem}

\begin{proof}
  \leanok
  The first statement is Lemma~\ref{lem:bundle-cov-deriv} between the preferred chart and the
  chart $j$; the second is the local curvature formula, as for Theorem~\ref{thm:conn-curvature-tensorial}.
\end{proof}

\begin{definition}[Bundle metrics and metric connections]
  \label{def:bundle-metric}
  \lean{LocallyConvex.Manifolds.IsLocalBundleMetric, LocallyConvex.Manifolds.metricD,
    LocallyConvex.Manifolds.IsLocalCompatible, LocallyConvex.Manifolds.BundleMetric,
    LocallyConvex.Manifolds.BundleMetric.IsPositive,
    LocallyConvex.Manifolds.BundleConnection.IsMetric}
  \leanok
  \uses{def:bundle-connection}
  A \emph{bundle metric} is a smooth family $h_i(y)$ of symmetric bilinear forms on $F$ in each
  trivialization, with $h_j(\tau y)(g\xi, g\eta) = h_i(y)(\xi, \eta)$; positivity is not
  required. A connection is \emph{metric} if
  $\partial_a h(\xi, \eta) = -h(\Gamma(a, \xi), \eta) - h(\xi, \Gamma(a, \eta))$ in every
  trivialization.
\end{definition}

For a weak metric a metric connection need not exist; the statements below assume one.

\begin{theorem}[Metric connections]
  \label{thm:metric-connections}
  \lean{LocallyConvex.Manifolds.IsLocalBundleMetric.dMB_apply,
    LocallyConvex.Manifolds.IsLocalBundleMetric.localCompat_sections,
    LocallyConvex.Manifolds.IsLocalBundleMetric.bundleCurv_skew,
    LocallyConvex.Manifolds.BundleConnection.IsMetric.curvLoc_skew,
    LocallyConvex.Manifolds.BundleConnection.IsMetric.curvature_skew}
  \leanok
  \uses{def:bundle-metric, thm:bundle-cov-deriv-global}
  For a metric connection, $X h(s, t) = h(\nabla_X s, t) + h(s, \nabla_X t)$ for $C^1$
  sections, and the curvature is skew-adjoint:
  $h(R(X, Y)s, t) + h(s, R(X, Y)t) = 0$.
\end{theorem}

\begin{proof}
  \leanok
  The first statement is the product rule for $h$. For the second, differentiate the
  compatibility condition in a second direction; the symmetry of the second derivatives of
  $y \mapsto h(y)(\xi, \eta)$ gives the identity, after the terms
  $h(\Gamma(a, \xi), \Gamma(b, \eta)) + h(\Gamma(b, \xi), \Gamma(a, \eta))$ cancel by symmetry.
\end{proof}

\begin{theorem}[Pullback connections]
  \label{thm:pullback-connection}
  \lean{LocallyConvex.Manifolds.TrivializationCover.pullbackRep,
    LocallyConvex.Manifolds.TrivializationCover.smoothMBOn_pullbackRep,
    LocallyConvex.Manifolds.TrivializationCover.pullbackRep_trans,
    LocallyConvex.Manifolds.BundleConnection.pullbackSymbol,
    LocallyConvex.Manifolds.BundleConnection.isLocalSymbol_pullbackSymbol,
    LocallyConvex.Manifolds.BundleConnection.pullback}
  \leanok
  \uses{def:pullback-bundle, def:bundle-connection}
  Let $f \colon N \to M$ be smooth and $\nabla$ a connection on $V$. In the trivialization of
  $f^*V$ over $k.\mathrm{source} \cap f^{-1}(U_i)$, let $F = \varphi_i \circ f \circ k^{-1}$. The
  symbols $\Gamma_i(F y)(\dd F(y) a, \xi)$ define a connection $f^*\nabla$ on $f^*V$.
\end{theorem}

\begin{proof}
  \leanok
  The local representatives intertwine the coordinate changes, $F_q \circ \tau' = \tau \circ F_p$,
  so by the chain rule the gauge law of $f^*\nabla$ is the gauge law of $\nabla$ at $F_p(y)$, in
  the direction $\dd F_p(y) a$.
\end{proof}

\begin{theorem}[Connections on Hom bundles]
  \label{thm:hom-connection}
  \lean{LocallyConvex.Manifolds.IsLocalSymbol.continuous_right,
    LocallyConvex.Manifolds.BundleConnection.homSymbolCLM,
    LocallyConvex.Manifolds.BundleConnection.homSymbol,
    LocallyConvex.Manifolds.BundleConnection.homSymbol_apply,
    LocallyConvex.Manifolds.BundleConnection.isLocalSymbol_homSymbol,
    LocallyConvex.Manifolds.BundleConnection.hom_transition,
    LocallyConvex.Manifolds.BundleConnection.hom}
  \leanok
  \uses{thm:hom-bundle, def:bundle-connection}
  Let $\nabla^1$, $\nabla^2$ be connections on vector bundles $V_1$, $V_2$ with normed fibers.
  The symbols $\Gamma^{\mathrm{Hom}}(y)(a, A)\xi = \Gamma_2(y)(a, A\xi) - A\,\Gamma_1(y)(a, \xi)$
  define a connection on $\mathrm{Hom}(V_1, V_2)$, the local form of
  $(\nabla_a A)\xi = \nabla^2_a(A\xi) - A(\nabla^1_a \xi)$.
\end{theorem}

\begin{proof}
  \leanok
  The map $y \mapsto \Gamma^{\mathrm{Hom}}(y)$ is smooth for the operator norm because
  $((y, a, A), \xi) \mapsto \Gamma^{\mathrm{Hom}}(y)(a, A)\xi$ is smooth (\cite{LCA}, theorem on
  operator-valued maps). For the gauge law, write the coordinate changes as
  $A \mapsto g_2 A g_1^{-1}$; differentiating $g_1(g_1^{-1}\xi) = \xi$ gives
  $\partial(g_1^{-1}) = -g_1^{-1}(\partial g_1)g_1^{-1}$, and the gauge laws of $\nabla^1$ and
  $\nabla^2$ give that of $\Gamma^{\mathrm{Hom}}$.
\end{proof}

\begin{corollary}[The dual connection]
  \label{cor:dual-connection}
  \lean{LocallyConvex.Manifolds.BundleConnection.trivial,
    LocallyConvex.Manifolds.BundleConnection.dual,
    LocallyConvex.Manifolds.BundleConnection.dual_Γ_apply}
  \leanok
  \uses{thm:hom-connection}
  A connection on a bundle $V$ with normed fiber induces a connection on the dual bundle $V'$,
  with symbols $\Gamma'(y)(a, \alpha)\xi = -\alpha(\Gamma(y)(a, \xi))$.
\end{corollary}

\begin{proof}
  \leanok
  Apply Theorem~\ref{thm:hom-connection} to $V$ and the trivial bundle $M \times \R$ with the
  trivial connection, whose symbols vanish.
\end{proof}

\begin{definition}[Refinements]
  \label{def:cover-refinement}
  \lean{LocallyConvex.Manifolds.TrivializationCover.Refines,
    LocallyConvex.Manifolds.TrivializationCover.refineRep,
    LocallyConvex.Manifolds.MBVectorBundleCore.refine,
    LocallyConvex.Manifolds.TrivializationCover.inf,
    LocallyConvex.Manifolds.TrivializationCover.infLeft,
    LocallyConvex.Manifolds.TrivializationCover.infRight}
  \leanok
  \uses{def:vector-bundle}
  A trivializing cover $(W_w)$ \emph{refines} $(V_k)$ along a map of indices $r$ if
  $W_w \subseteq V_{r(w)}$. A vector bundle trivialized over $(V_k)$ is trivialized over the
  refinement, with coordinate changes $g_{r(w) r(w')}$. Two covers $(U_i)$, $(V_k)$ have the
  common refinement $(U_i \cap V_k)$.
\end{definition}

\begin{theorem}[Connections on refinements; Hom bundles over different covers]
  \label{thm:refine-connection}
  \lean{LocallyConvex.Manifolds.BundleConnection.refineSymbol,
    LocallyConvex.Manifolds.BundleConnection.isLocalSymbol_refineSymbol,
    LocallyConvex.Manifolds.BundleConnection.refine,
    LocallyConvex.Manifolds.homCoreInf, LocallyConvex.Manifolds.BundleConnection.homInf}
  \leanok
  \uses{def:cover-refinement, thm:hom-connection}
  A connection on $V$ is written in the trivializations of a refinement with the symbols
  $\Gamma_{r(w)}(F y)(\dd F(y) a, \xi)$, where $F = \varphi_{r(w)} \circ \varphi_w^{-1}$. In
  particular, $\mathrm{Hom}(V_1, V_2)$ and its induced connection are defined for bundles
  trivialized over different covers, over their common refinement.
\end{theorem}

\begin{proof}
  \leanok
  As for pullback connections (Theorem~\ref{thm:pullback-connection}), with the change of charts
  $F$ in place of the local representative of $f$.
\end{proof}
